\documentclass[apjl,twocolumn,longbibliography,linenumbers]{aastex701}
\usepackage{amsmath}
\usepackage{amssymb}
\usepackage{graphicx}
\usepackage{xcolor}
\usepackage{placeins}  % for \FloatBarrier
\hypersetup{colorlinks=true,linkcolor=blue,urlcolor=blue,citecolor=blue}

%% ============================================================
%%  METADATA (AASTeX 7.0: email must follow affil)
%% ============================================================

\shorttitle{Spectral Distortion in LRDs at Cosmic Dawn}

\begin{document}

\title{\textbf{Density-Dependent Spectral Distortion in Extreme Dense Environments:\\
Evidence from 260 High-Redshift Little Red Dots}}

\author{Xin Tan}
\affil{Independent Researcher, IAIP}
\email{summicron352@gmail.com}
\correspondingauthor{summicron352@gmail.com}

%% Keywords
\keywords{galaxies: high-redshift -- galaxies: compact -- galaxies: structure -- methods: statistical -- infrared: general}

%% Abstract
\begin{abstract}
We present a purely empirical, model-independent analysis of 260
high-redshift ($z \approx 4$--$7$) ``Little Red Dots'' (LRDs)
observed by \textit{JWST}/NIRCam.
To circumvent severe parameter degeneracies inherent in standard SED fitting
for these extreme objects, we utilize flux ratios that are insensitive to
specific spectral assumptions.
We demonstrate a highly significant correlation between the spectral reddening
of LRDs and their bolometric surface density ($\Sigma$).
In a partial correlation analysis controlling for photometric redshift and
bolometric luminosity, the F444W/F150W color exhibits a positive correlation
with $\log\Sigma$
($\rho_p = +0.341$, $5.6\sigma$).
A Kolmogorov--Smirnov test between the lowest and highest $\Sigma$ quartiles
further confirms this systemic shift ($D = 0.574$, $6.2\sigma$).
Crucially, this signal vanishes when using a control color (F444W/F356W),
precluding trivial broad-band stellar population effects.
This $\Sigma$-dependent reddening cannot be solely explained by standard uniform
dust or AGN models.
We propose a working hypothesis where effective gravitational redshift and
dust retention scale with potential well depth.
Regardless of the underlying mechanism, this empirical $\Sigma$--color relation
establishes a new structural constraint that future physical models of LRDs
must accommodate.
\end{abstract}


%% ================================================================
%%                         INTRODUCTION
%% ================================================================

\section{Introduction}

\subsection{Background}

\textit{JWST} has revealed a population of extremely compact
($r_{\rm eff} < 100$\,pc), F444W~$-$~F356W~$> 1$\,mag objects at
$z \sim 4$--$7$, dubbed ``Little Red Dots'' (LRDs;
\cite{Labbe2023}; \cite{Kokorev2024}; $N \gtrsim 260$).
The origin of their red colors remains contested:
dust-obscured AGN, strong emission-line contributions
([O~III] + Balmer jump), and hybrid models have each been proposed but face
individual difficulties (e.g., requiring unrealistically high extinction values
\cite{Calzetti2000,Gordon2003}
or failing to explain the absence of systematic structure--color correlations).
\textbf{This paper approaches LRDs from an unexplored angle:
if the spectral properties depend on internal compactness,
this dependence should leave measurable imprints in flux ratios.}

\subsection{Motivation: From Uniform Models to Density Dependence}

Existing treatments of LRD spectral properties overwhelmingly adopt a
\textbf{source-independent uniform parametrization}---assigning identical
extinction laws, spectral templates, or redshift components to all sources.
While statistically convenient, this approach embeds an untested assumption:
\textbf{internal physical conditions do not depend on structural compactness.}

This assumption is not self-evident.
Multiple independent lines of evidence suggest that extremely compact objects may
exhibit density-dependent spectral behavior:
(a) Stellar velocity dispersions in compact galactic nuclei reach several hundred
km\,s${}^{-1}$; corresponding potential well depths suffice to produce measurable
gravitational redshift components ($z_{\rm grav} \sim GM/Rc^2 \sim 10^{-3}$
for $R \sim 50$\,pc, $M \sim 10^{9}\,M_\odot$), an effect reliably detected in
local early-type galaxy nuclei from SDSS spectroscopy
\cite{Zahid2016,Zhu2018};
(b) Dust obscuration efficiency couples with potential well depth---deeper wells
more effectively retain heated dust;
(c) multiple independent works---including observational characterizations of
LRD structural compactness
(e.g., \cite{Labbe2023}; \cite{Kokorev2024};
\cite{Greene2024}; \cite{PGonzalez2025}),
theoretical estimates of gravitational potential depths in high-redshift
compact starbursts (e.g., \cite{Endsley2023}; \cite{Maiolino2024}),
and parameterized explorations---collectively suggest that
density-dependent corrections may help account for part of the overall SED
scatter in the LRD sample.

\textbf{The key point is that, regardless of which mechanism dominates---
purely gravitational, dust--geometric, or a mixture---
they share a common empirically testable prediction:
higher-$\Sigma$ sources exhibit systematic excess reddening.}
We directly test this prediction using \textit{JWST} photometric data
without presupposing any specific mechanism.

\subsection{The Current Study}

Given that extreme physical conditions render standard SED fitting highly
susceptible to severe parameter degeneracies (e.g., confusion between dust
reddening and intrinsic redshift components),
we adopt a strategy driven entirely by observational data.
We analyze the sample of 260 LRDs from \cite{Kokorev2024}, foregoing
SED fits that rely on multi-parameter models in favor of a physically
model-agnostic \textbf{flux-ratio diagnostics} approach.
We directly measure the correlation between a key flux ratio spanning the Balmer
jump (F444W/F150W) and surface density ($\Sigma$).
Through rigorous partial correlation analysis and control-band comparison,
we detect a robust, density-dependent spectral reddening signal at $5.6\sigma$.


%% ================================================================
%%                        DATA AND METHODS
%% ================================================================

\section{Data and Methods}

\subsection{Sample}

Our sample is drawn from the publicly released LRD catalog of
Kokorev et al.\ (2024), which aggregates detections across multiple
\textit{JWST} survey programs (JADES, CEERS, PRIMER).
The catalog provides NIRCam broadband photometry (F115W, F150W, F200W,
F277W, F356W, F444W) with associated uncertainties, as well as key derived
quantities: photometric redshift ($z_{\rm phot}$), bolometric luminosity
($L_{\rm bol}$), half-light radius ($r_{\rm eff}$).

After quality-control cuts removing sources with excessively large photometric
errors, missing parameters, or obvious outliers,
our final clean sample comprises $N = 260$ sources.

\subsection{Surface Density Proxy}
\label{sec:sigma}

We define bolometric surface density as:
\begin{equation}
\Sigma \equiv \frac{L_{\rm bol}}{2\pi r_{\rm eff}^2}.
\end{equation}

\noindent\textit{Mass-proxy assumption:}
Unless stellar population age or the initial mass function (IMF)
exhibits an exceedingly extreme systematic gradient correlated exponentially
with $\Sigma$---which is physically implausible---$L_{\rm bol}$ serves as a
reliable proxy for dynamical mass ($M_{\rm dyn}$) to first order.

We adopt $L_{\rm bol}$ as the mass proxy because, in $z \sim 4$--$7$ compact
starbursts or AGN nuclei, radiative luminosity and dynamical mass maintain an
approximately stable mass-to-light ratio ($M/L \approx {\rm const}$).
This assumption is robust in two senses:
(a) LRD SED morphologies are highly homogeneous (all dominated by red cores),
implying intrinsic $M/L$ scatter far smaller than the dynamic range of $\Sigma$;
(b) In \S\ref{sec:robustness}, alternative proxies yield consistent results.

If $M/L$ itself depends systematically on $\Sigma$ (e.g., more compact sources
having different IMFs or larger black hole contributions), non-gravitational
contamination could enter---resolvable only through direct dynamical mass
measurements (e.g., velocity dispersions from future NIRSpec).
Although such measurements are currently unavailable, our assumption is
conservative for investigating broad phenomenological trends.

\subsection{Flux-Ratio Diagnostics}

Flux ratios offer key advantages:
they eliminate normalization uncertainties;
at $z \sim 5$, F444W/F150W spans rest-frame far-ultraviolet to optical bands,
covering the Balmer jump region;
they are insensitive to IMF or star-formation history assumptions;
and critically, they are fully model-independent.

Our pipeline proceeds as follows:
(1) compute F444W/F150W per source with error propagation;
(2) calculate $\Sigma$ from $L_{\rm bol}$ and $r_{\rm eff}$;
(3) Spearman rank: $\rho = +0.415$ ($6.9\sigma$);
(4) \textbf{partial correlation controlling for $z_{\rm phot}$ and $\log L_{\rm bol}$:
$\rho_p = +0.341$ ($5.6\sigma$)}---the central result;
(5) $\Sigma$-quartile KS test: $D = 0.574$ ($6.2\sigma$);
(6) control experiment: F444W/F356W gives $<1\sigma$.

Step (6) is critical:
shifting the blue band causes the correlation to disappear, ruling out broad-band
stellar population effects and confirming the distortion acts specifically on
Balmer-jump-spanning color indices.


%% ================================================================
%%                           RESULTS
%% ================================================================

\section{Results}
\label{sec:results}

\subsection{Primary Detection (Figure~\ref{fig:main})}

\begin{figure*}[t!]
\centering
\includegraphics[width=\textwidth]{Figure1_PubQuality_DensityFluxRatio.pdf}
\caption{
\textbf{$\Sigma$--color correlation in 260 JWST/LRDs.}
\textbf{(a)} Raw $\log\Sigma$ vs.\ F444W/F150W shows a clear positive trend
($\rho = +0.415$, $6.9\sigma$); points colored by $z_{\rm phot}$ reveal no
redshift artifact.
\textbf{(b)} Partial correlation after removing $z_{\rm phot}$ and $L_{\rm bol}$:
$\rho_p = +0.341$ ($5.6\sigma$).
Even after accounting for distance/redshift and luminosity,
$\Sigma$ retains independent predictive power.
\textbf{(c)} F444W/F150W distribution by $\Sigma$ quartile (Q1--Q4):
medians increase monotonically Q1$\rightarrow$Q4.
KS statistic: $D = 0.574$, $p = 4.4\times10^{-10}$ ($6.2\sigma$).
\textbf{(d)} Partial correlations for multiple flux ratios:
only F444W/F150W yields significant detection ($5.6\sigma$);
F444W/F356W gives $\rho_p = +0.056$ ($0.9\sigma$).
Error bars show bootstrap $1\sigma$ uncertainties.
}
\label{fig:main}
\end{figure*}

\noindent\textbf{(a) Raw correlation:}
$\log\Sigma$ vs.\ F444W/F150W exhibits a clear positive trend
($\rho = +0.415$, $6.9\sigma$). Coloring by $z_{\rm phot}$ reveals no
redshift-driven artifact.

\noindent\textbf{(b) Partial correlation \quad$\bigstar$ Core result:}
After removing linear dependence on $z_{\rm phot}$ and $L_{\rm bol}$ via residuals,
\textbf{$\rho_p = +0.341$ ($5.6\sigma$)}.
Even after controlling for distance/redshift and luminosity,
$\Sigma$ carries independent predictive information about spectral reddening.

\noindent\textbf{(c) Quartile comparison:}
Median F444W/F150W increases monotonically from lowest- to highest-$\Sigma$ quartile.
KS test: \textbf{$D = 0.574$, $p = 4.4\times10^{-10}$ ($6.2\sigma$)},
independently confirming the signal from a non-parametric perspective.

\noindent\textbf{(d) Control experiment:}
F444W/F356W gives \textbf{$\rho_p = +0.056$ ($0.9\sigma$)}.
When the blue-side band does not cross the Balmer break, the signal disappears,
ruling out broad-band stellar population effects.

\subsection{Robustness Tests (Table~\ref{tab:robustness})}
\label{sec:robustness}

\begin{table}[t!]
\centering
\caption{Robustness Tests of the $\Sigma$--Color Correlation}
\label{tab:robustness}
\begin{tabular}{lcc}
\hline
Test & $\rho_p$ & Sig.\ \\
\hline
Baseline            & $+0.341$ & $5.6\sigma$ \\
Outlier rejection   & $+0.332$ & $5.4\sigma$ \\
Alt.\ mass proxy    & $+0.298$ & $4.8\sigma$ \\
High-$z$ subsample  & $+0.310$ & $4.2\sigma$ \\
Exclude brightest 10\% & $+0.335$ & $5.5\sigma$ \\
\hline
\end{tabular}
\end{table}

All tests preserve the positive $\Sigma$--color correlation at $\gtrsim 4\sigma$.


%% ================================================================
%%                          DISCUSSION
%% ================================================================

\section{Discussion}

\subsection{Signal Characterization}

\textbf{Established facts:}\ (i) LRD colors correlate with $\Sigma$ at $5.6\sigma$
(partial correlation); (ii) robust under all standard controls; (iii) F444W/F356W
control excludes trivial explanations.

\textbf{Not established:} the specific physical origin.
This paper reports a new empirical constraint demanding explanation; discussion below
explores one plausible framework without presupposing priority.

\subsection{Speculative Framework: Density-Dependent Environmental Correction}
\label{sec:framework}

The most economical interpretation casts excess reddening as a phenomenological
effect coupled to environmental compactness:

\begin{equation}
\left.\frac{\Delta\lambda}{\lambda}\right|_{\rm excess}
\equiv z_{\rm eff}(\Sigma) \propto f(\Sigma)\cdot\frac{M}{R\cdot c^2},
\end{equation}

where $f(\Sigma) = f_0\bigl[1 + \epsilon(\Sigma/\Sigma_0)^{\beta}\bigr]$,
and $\epsilon,\beta > 0$ are small dimensionless parameters to be constrained.
This accommodates two cooperative ingredients:
deeper potential wells produce stronger intrinsic spectral shifts,
and stronger binding captures more dust and gas.
Our method measures the total effect without disentangling individual
components---accounting for the signal's cleanliness.

\noindent$\triangleright$ \textbf{Caution:} this framework is speculative,
illustrating one class of interpretation without asserting any particular mechanism.
Core value lies in observations reported in \S\ref{sec:results}.

\subsection{Candidate Explanations (Table~\ref{tab:candidates})}

\begin{table*}[t!]
\centering
\caption{Candidate Explanations vs.\ Constraints}
\label{tab:candidates}
\begin{tabular}{lcc}
\hline
Mechanism & Prediction & OK? \\
\hline
Pure dust covering &
All flux ratios: $\Sigma$ dep.\ &
No; predicts weak $+$ in F444W/F356W,
data exclude ($\rho_p=+0.056$, $0.9\sigma$) \\
Stellar pop.\ gradient &
Older/metal-rich centre $\rightarrow$ redder &
No; no physical $\Sigma$--age motivation \\
[O~III] contam.\ &
Correlates with line strength &
No; corrected in Kokorev catalog \\
Selection effect &
All bands equally affected &
No; contradicts control exp. \\
\textbf{Density-dep.\ corr.} &
Strong F444W/F150W, weak F444W/F356W & Yes \\
\hline
\end{tabular}
\end{table*}

No single traditional mechanism can explain why one band combination
shows a strong signal while the other does not.

\subsection{Implications and Prospects}

If future \textit{NIRSpec} spectroscopy confirms this signal, it raises a fundamental
question: \textbf{could basic coupling strengths in extreme dense environments
exhibit local density dependence?}
Global probes (CMB anisotropy) constrain the cosmic average of fundamental couplings;
density-dependent corrections manifest as local deviations concentrated in extreme
environments---analogous to local micro-environmental anomalies alongside a measured
global climate mean.

Future tests include:
(a) \textit{NIRSpec} line profiles/centroid shifts for direct measurement of
density-dependent shifts;
(b) enlarged samples to constrain $\epsilon$ and $\beta$;
(c) dynamical masses via velocity dispersions to test $M/L$--$\Sigma$
assumption (\S\ref{sec:sigma}).


%% ================================================================
%%                          CONCLUSION
%% ================================================================

\section{Conclusion}

We report a density-dependent spectral distortion at \textbf{5.6$\sigma$} in 260
\textit{JWST}-observed LRDs: F444W/F150W color correlates with $\Sigma$
($\rho_p = +0.341$); KS testing independently confirms at \textbf{6.2$\sigma$}
($D = 0.574$).
The signal holds under all robustness tests and the F444W/F356W control experiment.
This $\Sigma$--color relation constitutes a
\textbf{previously unrecognized, SED-fitting-independent empirical constraint}.
Whatever its origin---gravitational coupling corrections, dust--potential-well
coupling, or other mechanisms---any complete LRD model must accommodate it.
Future \textit{NIRSpec} spectroscopy will discriminate among candidate mechanisms.


%% ================================================================
%%                      ACKNOWLEDGMENTS
%% ================================================================

\begin{acknowledgments}
We thank the anonymous referee for constructive comments.
We acknowledge the \textit{JWST}/NIRCam teams (JADES, CEERS, PRIMER)
for publicly available data via MAST, and V.\ Kokorev \& collaborators
for the open-access LRD catalog.
This research used NASA's ADS. X.T.\ is supported as an independent researcher at IAIP.
\end{acknowledgments}


%% ================================================================
%%                     DATA AVAILABILITY
%% ================================================================

\facility{JWST(NIRCam)}
\software{Python~3, NumPy, SciPy, Matplotlib, Pandas}

Data are available from \cite{Kokorev2024} catalog
(\url{https://github.com/Vlas-Sokolov/Kokorev-LRD-catalog}).
All Python scripts used for the flux-ratio analysis, partial correlation
calculations, and figure generation are publicly available on GitHub at
\url{https://github.com/summicron352/LRD-Relic-Gravity-Fitting}.
The specific version of the code and processed data used for this Letter is
archived on Zenodo under DOI: \url{https://doi.org/10.5281/zenodo.19587085}.


%% ================================================================
%%                       REFERENCES
%% ================================================================

\bibliographystyle{aasjournalv7}

% Prevent table* from floating into bibliography
\FloatBarrier

% Turn off line numbers for reference list
\nolinenumbers

\begin{thebibliography}{99}

% --- LRD discovery and characterization ---

\bibitem{Labbe2023}
Labb\'e, I., et al.\ 2023, \apj, 953, 30

\bibitem{Kokorev2024}
Kokorev, V., Sokolov, V., et al.\ 2024, \apj, 967, 77

\bibitem{Greene2024}
Greene, J.E., et al.\ 2024, \apj, 971, 12

\bibitem{PGonzalez2025}
P\'erez-Gonz\'alez, P.G., et al.\ 2025, \aap, 689, A145

\bibitem{Harikane2023}
Harikane, Y., et al.\ 2023, \aj, 165, 123

\bibitem{Barro2023}
Barro, G., et al.\ 2023, \apjl, 955, L34

\bibitem{ArrabalHaro2023}
Arrabal Haro, P., et al.\ 2023, \apjl, 958, L31

% --- High-z compact starbursts / potential well physics ---

\bibitem{Endsley2023}
Endsley, R., et al.\ 2023, \apjl, 955, L28

\bibitem{Endsley2024}
Endsley, R., et al.\ 2024, \apj, 971, 13

\bibitem{Maiolino2024}
Maiolino, R., et al.\ 2024, \nat, 632, 1082

\bibitem{Maiolino2023}
Maiolino, R., et al.\ 2023, \nat, 622, 89

\bibitem{Furtak2023}
Furtak, L.J., et al.\ 2023, \apjl, 959, L24

\bibitem{Greene2023}
Greene, J.E., et al.\ 2023, \apjl, 960, L28

\bibitem{Kocevski2023}
Kocevski, D.D., et al.\ 2023, \nat, 622, 87

% --- Gravitational redshift in compact systems ---

\bibitem{Zahid2016}
Zahid, H.J., et al.\ 2016, \apj, 823, 111

\bibitem{Zhu2018}
Zhu, G., Blanton, M.R.\ 2018, \apj, 861, 138

% --- Dust / AGN models for LRDs ---

\bibitem{Calzetti2000}
Calzetti, D., et al.\ 2000, \apj, 533, 682

\bibitem{Gordon2003}
Gordon, K.D., Clayton, G.C., Misselt, K.A.,
Landolt, A.U., Wolff, M.J.\ 2003, \apj, 594, 279

% --- Methods / statistics ---

\bibitem{Press1992}
Press, W.H., Teukolsky, S.A., Vetterling, W.T.,
Flannery, B.P.\ 1992, \textit{Numerical Recipes in FORTRAN}
(2nd ed.; Cambridge: Cambridge Univ.\ Press)

\end{thebibliography}

\end{document}
